\documentclass[12pt]{amsart}
\usepackage[utf8]{inputenc}
\usepackage[T1]{fontenc}
\usepackage{lmodern}
\usepackage{amsmath,amssymb,amsthm,mathtools}
\usepackage{microtype}
\usepackage[colorlinks=true,linkcolor=blue,citecolor=blue,urlcolor=blue]{hyperref}

\newtheorem{theorem}{Theorem}
\newtheorem{lemma}[theorem]{Lemma}
\newtheorem{proposition}[theorem]{Proposition}
\theoremstyle{remark}
\newtheorem*{remark}{Remark}
\newtheorem*{remarks}{Remarks}

\newcommand{\C}{\mathbb{C}}
\newcommand{\R}{\mathbb{R}}
\newcommand{\Edown}{E_{\downarrow}}
\newcommand{\Argdown}{\operatorname{Arg}_{\downarrow}}
\newcommand{\Logdown}{\operatorname{Log}_{\downarrow}}
\newcommand{\Logpr}{\operatorname{Log}}
\newcommand{\LE}{\mathcal{LE}}
\newcommand{\Term}{\mathcal{T}}

\title[Obstructions to self-seeding Sheffer operators]{Obstructions to Self-Seeding Sheffer Operators for Calculator-Style Elementary Functions}
\author{Marouane Lamharzi Alaoui}
\thanks{Firassa AI. Email address: \texttt{marouane@firassa.ai}.}
\date{April 13, 2026}

\begin{document}

\begin{abstract}
Odrzywo\l{}ek recently showed that the pair $(1,\operatorname{eml})$, with $\operatorname{eml}(x,y)=e^x-\log y$, generates a fixed scientific-calculator basis of elementary functions and constants.  We prove that the most natural self-seeding analytic strategy fails: any holomorphic or real-analytic binary germ with constant diagonal $f(x,x)\equiv c$ generates only term germs that fix $c$, and any purely real logarithmico-exponential primitive generates only eventually monotone unary tail-germs, hence never $\sin$ or $\cos$.  We also isolate the branch-correct core identities behind the EML reconstruction under a lower-edge logarithm convention, explain why the remaining published calculator-basis verification is mixed symbolic--numerical rather than a full symbolic theorem, and formulate multi-step self-seeding via off-diagonal fibers as the precise remaining open problem.
\end{abstract}

\maketitle

\section{Introduction}

Single reusable primitives often reveal hidden structure.  The classical example is the Sheffer stroke, which shows that Boolean logic can be generated from one binary connective \cite{Sheffer1913}.  Odrzywo\l{}ek has recently proposed a continuous analogue: the binary operator
\begin{equation}\label{eq:eml}
\operatorname{eml}(x,y)=e^x-\log y,
\end{equation}
and argued that, together with the constant $1$, it generates a concrete scientific-calculator basis of elementary functions, constants, and arithmetic operations \cite{Odrzywolek2026}.  His paper is explicit about scope: the claim concerns the finite calculator basis listed in Table~1 of \cite{Odrzywolek2026}, not the full Liouville--Ritt class of elementary extensions.  It nevertheless raises three natural questions:
\begin{itemize}
    \item can one replace the pair $(1,\operatorname{eml})$ by a single \emph{self-seeding} binary operator;
    \item can one state the EML reconstruction with a mathematically correct branch convention;
    \item and is a purely real tame Sheffer operator possible at all?
\end{itemize}

At the set-theoretic level, the first question is already settled.  Evans proved that every countable clone can be embedded into one generated by a single binary operation on a countable set \cite{Evans1989}, and Goldstern observed more generally that on any infinite set one binary operation suffices for any finite or countable family of target operations \cite{Goldstern2010}.  The infinite clone lattice is correspondingly huge and resists any Post-style global classification; see the survey of Goldstern and Pinsker \cite{GoldsternPinsker2008}.  Thus the mathematically interesting problem is not mere existence, but existence \emph{inside a tame regularity class}: continuous, real-analytic, holomorphic, or logarithmico-exponential.

The present note makes three theorem-level contributions and one clarifying one.

\smallskip
\noindent\textbf{(1) A local obstruction for one-step analytic self-seeding.}
If an analytic binary germ satisfies $f(x,x)\equiv c$, then every term germ generated by $f$ and the coordinate projections belongs to the maximal ideal $(z_1-c,\dots,z_n-c)$ up to the constant term $c$.  Equivalently, every unary term fixes $c$, and the only constant term operation is the constant $c$ itself.  This rules out the entire constant-diagonal analytic strategy.

\smallskip
\noindent\textbf{(2) A real obstruction for logarithmico-exponential primitives.}
If the primitive binary operator is a real logarithmico-exponential (LE) function, then every unary term in the resulting clone is again LE on some tail.  Such germs lie in a Hardy field and are eventually monotone.  Hence no nonconstant periodic function, in particular neither $\sin$ nor $\cos$, can occur.

\smallskip
\noindent\textbf{(3) A branch-correct core for the EML reconstruction.}
Odrzywo\l{}ek's text uses the principal logarithm in prose but later notes a sign mismatch on the negative real axis \cite{Odrzywolek2026}.  We isolate the clean fix: inside EML one should use the lower-edge logarithm, with argument in $[-\pi,\pi)$.  Under this convention the reconstructed logarithm becomes the usual principal logarithm, and the sign of $\log(-1)$ comes out correctly.  What we do \emph{not} claim is a new full theorem proving the entire Table~1 witness chain from scratch: the public verifier for the larger chain is mixed symbolic--numerical, and we describe its status honestly below.

\smallskip
\noindent\textbf{(4) A geometric formulation of the remaining open problem.}
The diagonal obstruction does \emph{not} settle the multi-step case: Odrzywo\l{}ek already noted the affine trap
\[
B(x,y)=x-\frac y2,
\qquad B(B(x,x),x)=0,
\]
which generates a constant only after a second step \cite{Odrzywolek2026}.  We show that any minimal-depth constant unary term of depth at least $2$ forces an off-diagonal analytic curve inside a level set $f^{-1}(k)$.  This identifies the correct geometric search space for any future tame positive construction.

These results are intentionally narrower than the full theory of elementary functions in the sense of differential algebra.  Ritt's and Risch's classical works remain the right ambient framework for Liouville towers, exponential and logarithmic adjunctions, and the constant problem \cite{Ritt1948,Risch1969}.  Our point is different: once the target is Odrzywo\l{}ek's concrete calculator basis, one can prove meaningful impossibility theorems for self-seeding binary primitives without solving the full Liouville--Ritt problem.

\section{Preliminaries}

We fix some notation.

\subsection*{Analytic germs and term germs}
Let $K\in\{\R,\C\}$.  For $p\in K^n$, let $\mathcal{A}_{n,p}(K)$ denote the ring of germs at $p$ of real-analytic maps $K^n\to K$ if $K=\R$, and of holomorphic maps $K^n\to K$ if $K=\C$.

Fix $c\in K$ and write $p_n=(c,\dots,c)\in K^n$.  Let $f$ be a germ at $(c,c)\in K^2$.  We define recursively the set $\Term_n(f;c)\subseteq \mathcal{A}_{n,p_n}(K)$ of $n$-ary \emph{term germs generated by $f$} as follows:
\begin{itemize}
    \item each coordinate projection $\pi_j(z_1,\dots,z_n)=z_j$ belongs to $\Term_n(f;c)$;
    \item if $u,v\in \Term_n(f;c)$ and $u(p_n)=v(p_n)=c$, then the composite germ $f(u,v)$ belongs to $\Term_n(f;c)$.
\end{itemize}
In all our local arguments the condition $u(p_n)=v(p_n)=c$ is proved inductively, so the compositions are well-defined after possibly shrinking neighborhoods.

For later use we define the depth of a unary term recursively by
\[
\operatorname{depth}(x)=0,
\qquad
\operatorname{depth}(f(u,v))=1+\max\{\operatorname{depth}(u),\operatorname{depth}(v)\}.
\]

\subsection*{Logarithmico-exponential functions and Hardy fields}
For each $m\geq 1$, let $\LE_m$ denote the class of partial real functions in $m$ variables obtained from real constants and the coordinate projections by finitely many uses of the field operations, the real exponential, the real logarithm on positive arguments, and composition; the domain of such an expression is defined recursively by the usual pullback rules.  This is the standard class of logarithmico-exponential functions in the sense of Hardy and later computer-algebra work; see Hardy, Rosenlicht, and Strzebo\'nski \cite{Hardy1910,Rosenlicht1983,Strzebonski2019}.

A \emph{Hardy field} is a field of germs at $+\infty$ of real-valued functions on tails $(a,\infty)$ that is closed under differentiation.  The germs of logarithmico-exponential functions defined near $+\infty$ form a Hardy field, and every element of a Hardy field is eventually monotone \cite{Rosenlicht1983,AschenbrennerVanDenDriesVanDerHoeven2017}.

\subsection*{The target basis}
We write $\mathcal{B}_{\mathrm{calc}}$ for Odrzywo\l{}ek's finite scientific-calculator basis from Table~1 of \cite{Odrzywolek2026}.  It contains the constants $\pi,e,i,-1,1,2$, the usual real variables, arithmetic operations, exponentials and logarithms, radicals and powers, and the trigonometric and hyperbolic families together with their inverses.  Every completeness statement in this paper refers to that concrete basis, not to the full Liouville--Ritt class.

\section{The diagonal ideal obstruction}\label{sec:diagonal}

We begin with the most natural constant-free strategy: look for an analytic binary operator $f$ with a constant diagonal,
\[
f(x,x)\equiv c.
\]
The next theorem shows that such an operator can never seed more than that one constant.

\begin{theorem}[Local diagonal ideal obstruction]\label{thm:diagonal}
Let $K\in\{\C,\R\}$, let $c\in K$, and let $f$ be a germ at $(c,c)\in K^2$ of a holomorphic map if $K=\C$, respectively a real-analytic map if $K=\R$.  Assume that
\[
f(x,x)\equiv c
\]
as a germ along the diagonal.

Then, for every $n\geq 1$ and every term germ $t\in \Term_n(f;c)$, there exist analytic germs
\[
h_1,\dots,h_n\in \mathcal{A}_{n,p_n}(K)
\]
at $p_n=(c,\dots,c)$ such that
\begin{equation}\label{eq:diagonal-ideal}
t(z_1,\dots,z_n)=c+\sum_{j=1}^n (z_j-c)h_j(z_1,\dots,z_n).
\end{equation}
Consequently every term germ satisfies $t(p_n)=c$; the only constant term operation obtainable from arbitrary input is the constant $c$ itself; and every unary term germ fixes $c$.  In particular, any target basis requiring either a second constant $d\neq c$ or a unary operation $u$ with $u(c)\neq c$ is unreachable.
\end{theorem}

\begin{proof}
The key local factorization is
\begin{equation}\label{eq:factorization}
f(x,y)=c+(x-y)g(x,y)
\end{equation}
for a suitable analytic germ $g$ at $(c,c)$.  Since the statement is local, we translate $c$ to $0$ and work with the germ
\[
F(X,Y):=f(c+X,c+Y)-c
\]
at $(0,0)$.  The hypothesis becomes $F(T,T)\equiv 0$.  Now set
\[
H(U,V):=F(U+V,V).
\]
Then $H$ is analytic and $H(0,V)=F(V,V)\equiv 0$.  Therefore the convergent power series of $H$ at $(0,0)$ has no $U^0$ term, so it is divisible by $U$:
\[
H(U,V)=U\,G(U,V)
\]
for some analytic germ $G$.  Returning to the original coordinates gives
\[
F(X,Y)=(X-Y)G(X-Y,Y),
\]
hence \eqref{eq:factorization} with $g(x,y)=G(x-y,y-c)$ after undoing the translation.

We now prove \eqref{eq:diagonal-ideal} by induction on the complexity of the term $t$.

For a projection $\pi_k(z_1,\dots,z_n)=z_k$ we have
\[
\pi_k(z_1,\dots,z_n)=c+(z_k-c),
\]
so \eqref{eq:diagonal-ideal} holds with $h_k\equiv 1$ and $h_j\equiv 0$ for $j\neq k$.

Suppose next that
\[
u(z)=c+\sum_{j=1}^n(z_j-c)a_j(z),
\qquad
v(z)=c+\sum_{j=1}^n(z_j-c)b_j(z)
\]
with analytic germs $a_j,b_j$ at $p_n$, and define $t=f(u,v)$.  By the induction hypothesis, $u(p_n)=v(p_n)=c$, so the composite germ is defined.  Using \eqref{eq:factorization},
\[
\begin{aligned}
t(z)
&=f(u(z),v(z))
 = c+\bigl(u(z)-v(z)\bigr)g\bigl(u(z),v(z)\bigr) \\
&= c+\sum_{j=1}^n (z_j-c)\bigl(a_j(z)-b_j(z)\bigr)g\bigl(u(z),v(z)\bigr),
\end{aligned}
\]
which is exactly of the required form.

This completes the induction.  The stated consequences follow immediately by substituting $z_1=\dots=z_n=c$ in \eqref{eq:diagonal-ideal}.\qedhere
\end{proof}

\begin{remarks}
\leavevmode
\begin{enumerate}
    \item The theorem is intentionally local.  No convexity hypothesis is needed.  The factorization argument takes place in the local analytic ring, so non-convex global domains and branch-cut geometry play no role.
    \item The same proof works verbatim for holomorphic and real-analytic germs.  Because the theorem is stated for germs at $(c,c)$, there is no boundary-value loophole: the seed point is necessarily an interior point of the local domain.
    \item The holomorphic or real-analytic hypothesis at the seed point is essential.  The theorem does \emph{not} extend to meromorphic operators singular at $(c,c)$.  Indeed,
    \[
    f(z,w)=\frac{z-w}{z+w}
    \]
    satisfies $f(x,x)=0$ for $x\neq 0$, yet
    \[
    f(f(x,x),x)=f(0,x)=-1
    \qquad (x\neq 0).
    \]
    Singularities at the seed point can therefore create escape routes unavailable in the analytic setting.
    \item One should not overread the theorem.  It does \emph{not} imply that specific unary targets such as $\exp$ or a chosen logarithm branch are individually impossible: a unary target that happens to fix $c$ is not excluded by the fixed-point argument alone.  The actual obstruction is different and stronger for the calculator problem: the generated clone can realize only one constant and all unary terms fix the same point $c$.  This already rules out Odrzywo\l{}ek's basis $\mathcal{B}_{\mathrm{calc}}$, which contains several distinct constants \cite{Odrzywolek2026}.
\end{enumerate}
\end{remarks}

Theorem~\ref{thm:diagonal} kills the most obvious ``binary EML without the external $1$'' strategy.  In particular, any analytic candidate satisfying $f(x,x)=1$ can at best reproduce the constant $1$; it cannot recover a second named constant such as $2$, $-1$, $e$, $i$, or $\pi$ from arbitrary input.

\section{The Hardy-field obstruction for real LE primitives}\label{sec:hardy}

The previous theorem treats analytic operators with a constant diagonal.  We now turn to the purely real problem that Odrzywo\l{}ek also singled out: can one hope for a tame real-domain Sheffer operator that generates the trigonometric family without passing through complex logarithmic branching?  For logarithmico-exponential primitives the answer is no.

\begin{theorem}[Hardy-field obstruction for real LE primitives]\label{thm:hardy}
Let $f\in \LE_2$, and let $t$ be a unary term obtained from $f$ and the identity projection by finite composition.  Assume that $t$ is defined on some tail $(a,\infty)$.

Then there exists $b\geq a$ such that $t\in \LE_1$ on $(b,\infty)$.  In particular, the germ of $t$ at $+\infty$ belongs to a Hardy field, hence is eventually monotone.  Consequently no nonconstant periodic function, and more generally no function that fails to be eventually monotone, can lie in the unary term clone generated by $f$.

Therefore no purely real binary logarithmico-exponential primitive can generate Odrzywo\l{}ek's calculator basis $\mathcal{B}_{\mathrm{calc}}$, because that basis contains $\sin$ and $\cos$.
\end{theorem}

\begin{proof}
We first record the closure under substitution carefully.  By definition, the class $\LE_m$ is built from constants and coordinate projections by the field operations, $\exp$, $\log$ on positive arguments, and composition.  If $g,h\in \LE_1$ on some tails, then $g\pm h$, $gh$, and $g/h$ are again logarithmico-exponential on the subdomain where they are defined.  Likewise $\exp(g)$ is logarithmico-exponential wherever $g$ is defined, and $\log(g)$ is logarithmico-exponential on the subdomain where $g>0$.  Finally, if $F\in \LE_2$ and $u,v\in \LE_1$, then $x\mapsto F(u(x),v(x))$ is again an LE expression on the pullback domain
\[
\{x:(u(x),v(x))\in \operatorname{dom}(F)\}.
\]
Thus the class is closed under substitution, with domains tracked by pullback.

Now argue by induction on the term depth.  The identity projection belongs to $\LE_1$.  If a term $t$ has the form $t(x)=f(u(x),v(x))$, and if $u$ and $v$ are logarithmico-exponential on some tails, then by the previous paragraph so is $t$ on the pullback domain.  Since by hypothesis $t$ is actually defined on some tail $(a,\infty)$, we may restrict to a smaller tail $(b,\infty)$ contained in that domain.  The induction proves that every unary term defined near $+\infty$ belongs to $\LE_1$ near $+\infty$.

Classical Hardy-field theory shows that logarithmico-exponential germs at $+\infty$ form a Hardy field \cite{Hardy1910,Rosenlicht1983,AschenbrennerVanDenDriesVanDerHoeven2017}.  Let $[t]$ be the germ of our unary term.  Since the Hardy field is closed under differentiation, $[t']$ also belongs to it.  If $[t']=0$, then $t$ is eventually constant.  If $[t']\neq 0$, then every representative of $[t']$ has eventually constant sign; hence $t$ is eventually strictly monotone.

Therefore every unary term in the clone is eventually monotone.  Any function that fails to be eventually monotone is excluded.  In particular no nonconstant periodic function can occur, because a periodic function oscillates on every tail.  Since Odrzywo\l{}ek's basis contains $\sin$ and $\cos$, no purely real $f\in \LE_2$ can generate the full basis.\qedhere
\end{proof}

\begin{remark}
Theorem~\ref{thm:hardy} excludes more than periodic functions: it rules out every unary target that is not eventually monotone, including oscillations of growing amplitude.  At the same time its hypothesis is genuinely narrower than ``real analytic.''  The theorem does \emph{not} rule out an arbitrary real-analytic primitive outside the LE class.  The gap between ``real LE'' and ``real analytic'' is exactly where a positive tame real construction, if one exists, would have to hide.
\end{remark}

\section{Branch-correct EML identities and the status of the published verification}\label{sec:eml}

Odrzywo\l{}ek's EML reconstruction is branch-sensitive.  The prose of \cite{Odrzywolek2026} speaks of the principal logarithm, but the paper also notes a sign mismatch on the negative real axis and suggests either redefining the branch inside EML or manually correcting the sign of $i$.  The public verification script in the accompanying repository takes the first route: it defines a lower-edge logarithm inside EML, with argument in $[-\pi,\pi)$ \cite{OdrzywolekVerifier2026}.  We make that convention explicit and prove the core identities symbolically.

Define
\[
\Argdown:\C^{\times}\to[-\pi,\pi),
\qquad
\Logdown z:=\log|z|+i\Argdown(z),
\]
and
\begin{equation}\label{eq:Emldown}
\Edown(x,y):=e^x-\Logdown(y).
\end{equation}
Thus $\Logdown(x)=\log|x|-i\pi$ for negative real $x$.

\begin{proposition}[Branch-correct core EML identities]\label{prop:emlcore}
Let $\Edown$ be defined by \eqref{eq:Emldown}, and define
\[
\exp_E(z):=\Edown(z,1),
\qquad
\log_E(z):=\Edown\bigl(1,\exp_E(\Edown(1,z))\bigr)
\qquad (z\in\C^{\times}).
\]
Then
\[
\exp_E(z)=e^z
\qquad (z\in\C),
\]
and
\[
\log_E(z)=\Logpr(z)
\qquad (z\in\C^{\times}),
\]
where $\Logpr$ denotes the usual principal logarithm with argument in $(-\pi,\pi]$.  In particular,
\[
\log_E(-1)=i\pi.
\]
Moreover,
\[
0=\Edown\bigl(1,\Edown(\Edown(1,1),1)\bigr)
\]
as a symbolic identity.
\end{proposition}

\begin{proof}
Since $\Logdown(1)=0$, we have immediately
\begin{equation}\label{eq:exp-witness}
\exp_E(z)=\Edown(z,1)=e^z.
\end{equation}

Now write $z=re^{i\theta}$ with $r>0$ and $\theta\in(-\pi,\pi]$.  By definition,
\[
\Edown(1,z)=e-\Logdown(z),
\qquad
\exp_E\bigl(\Edown(1,z)\bigr)=\exp\bigl(e-\Logdown(z)\bigr)=\frac{e^e}{z}.
\]
The point $e^e/z=(e^e/r)e^{-i\theta}$ has lower-edge argument $-\theta$, which indeed lies in $[-\pi,\pi)$ because $\theta\in(-\pi,\pi]$.  Therefore
\[
\Logdown\!\left(\frac{e^e}{z}\right)=e-\log r-i\theta.
\]
Substituting into the definition of $\log_E$ gives
\[
\log_E(z)=e-\left(e-\log r-i\theta\right)=\log r+i\theta=\Logpr(z),
\]
which proves the branch-correct logarithm identity.  In particular,
\begin{equation}\label{eq:branch-sign}
\log_E(-1)=i\pi.
\end{equation}

Finally,
\[
\Edown\bigl(1,\Edown(\Edown(1,1),1)\bigr)
= e-\Logdown(e^e)=e-e=0,
\]
so the displayed zero witness is a symbolic identity as claimed.\qedhere
\end{proof}

\begin{remark}[What is and is not proved here]\label{rem:status}
Proposition~\ref{prop:emlcore} is the theorem-level part of the EML section.  It proves the correct branch convention and the key negative-axis sign identity \eqref{eq:branch-sign}.  The public Mathematica verifier in Odrzywo\l{}ek's repository then checks a much larger finite witness chain for the constants, arithmetic operations, and the trigonometric and hyperbolic families \cite{OdrzywolekVerifier2026}.  However, that public verification is mixed symbolic--numerical rather than purely symbolic: the script comments out direct checks for \texttt{logEML}, routes several later identities through so-called phase-2 canonical witnesses, marks some earlier witnesses as ``old flaky branch/sign witnesses,'' and allows a unary identity to pass via a finite list of transcendental sample points if \texttt{FullSimplify} does not discharge it.  For that reason we do \emph{not} elevate the full Table~1 reconstruction to theorem status in the present note.

The honest conclusion is therefore the following.  The lower-edge convention above resolves the branch problem mathematically and yields the correct symbolic core.  Beyond that core, the currently public evidence for the full Table~1 reconstruction is a branch-correct \emph{computer-assisted verification claim}, not a completely symbolic proof from first principles.  This does not affect Theorems~\ref{thm:diagonal} and~\ref{thm:hardy}, which are independent of the external verifier.
\end{remark}

\section{Discussion and the remaining open problem}\label{sec:discussion}

Theorem~\ref{thm:diagonal} rules out the entire one-step constant-diagonal analytic family, but it does \emph{not} settle multi-step self-seeding.  Odrzywo\l{}ek emphasized exactly this point with the affine example
\[
B(x,y)=x-\frac y2,
\qquad B(B(x,x),x)=0,
\]
which seeds a constant at depth $2$ even though $B(x,x)=x/2$ is not constant \cite{Odrzywolek2026}.  The next lemma isolates the geometry that any such multi-step construction must exploit.

\begin{lemma}[Off-diagonal fiber criterion]\label{lem:fibers}
Let $K\in\{\R,\C\}$, let $f$ be real-analytic if $K=\R$ or holomorphic if $K=\C$, and let
\[
t(x)\equiv k
\]
be a constant unary term of minimal depth $m\geq 2$ in the term language generated by $f$ and the identity projection.  Write
\[
t(x)=f(u(x),v(x))
\]
with $u$ and $v$ of smaller depth.

Then $u$ and $v$ are nonconstant, $u-v$ is not identically zero, and the analytic curve
\[
x\longmapsto (u(x),v(x))
\]
lies in the level set $f^{-1}(k)$ and is not contained in the diagonal.  If, in addition, the image curve meets a regular point of $f^{-1}(k)$, then near that point one of $u$ and $v$ is a local analytic function of the other.
\end{lemma}

\begin{proof}
Since $t=f(u,v)$ is a constant unary term of minimal depth $m\geq 2$, neither $u$ nor $v$ can be constant: otherwise there would already exist a constant unary term of smaller depth.

We next show that $u-v$ cannot vanish identically.  Suppose, for contradiction, that $u\equiv v$.  Then
\[
t(x)=f(u(x),u(x)).
\]
Define the diagonal unary germ
\[
D(z):=f(z,z).
\]
We then have
\[
D(u(x))\equiv k.
\]
Because $u$ is nonconstant analytic, this forces $D$ itself to be the constant germ $k$.  In the holomorphic case this follows from the open mapping theorem: the image of a nonconstant holomorphic germ contains an open neighborhood, so $D$ is constant on an open set and hence everywhere as a germ by the identity theorem.  In the real-analytic case, $u'$ is not identically zero, so there exists $x_0$ with $u'(x_0)\neq 0$; by the inverse function theorem, $u$ maps a neighborhood of $x_0$ onto an open interval, and $D$ is constant on that interval, hence constant as a real-analytic germ.  Thus the depth-$1$ unary term $x\mapsto f(x,x)$ is already constant, contradicting the minimality of $m\geq 2$.  Therefore $u-v$ is not identically zero.

Since $t(x)\equiv k$, we have
\[
f(u(x),v(x))=k
\]
for all $x$ in the domain of the term, so the image curve lies in the fiber $f^{-1}(k)$.  The previous paragraph shows that the curve is not contained in the diagonal.

Finally, if the image curve meets a regular point $(a,b)$ of the level set $f^{-1}(k)$, then by the implicit function theorem that level set is locally a one-dimensional analytic graph near $(a,b)$; hence locally one of the coordinates is an analytic function of the other.\qedhere
\end{proof}

As an immediate consequence, if every fiber $f^{-1}(k)$ contains no nonconstant analytic curve off the diagonal, then no constant unary term of depth at least $2$ can exist.

Lemma~\ref{lem:fibers} is not a completeness theorem.  It is a geometric description of what the remaining positive problem would have to look like.  Any tame multi-step self-seeding candidate must have a \emph{nonconstant} diagonal and, at the same time, sufficiently rich off-diagonal constant fibers to carry analytic curves of the kind described in the lemma.  In particular, one should not search among operators whose only mechanism for generating a constant is diagonal cancellation: Theorem~\ref{thm:diagonal} shows that this entire family is impossible.

The unrestricted clone-theoretic constructions point in the same qualitative direction.  Classical one-binary-generator constructions on infinite sets do not seed constants by diagonal cancellation; rather, they use a nonconstant diagonal to create new encodings or successor states.  See Evans and Goldstern \cite{Evans1989,Goldstern2010}.  This strongly suggests that if a tame positive solution exists, it will have to use the diagonal productively rather than cancellatively.

Two concrete open problems emerge.
\begin{itemize}
    \item \textbf{Tame multi-step self-seeding.}  Does there exist a holomorphic or real-analytic binary germ with nonconstant diagonal whose term clone nevertheless contains a nontrivial constant unary term?
    \item \textbf{The real-analytic gap.}  Theorem~\ref{thm:hardy} rules out all purely real LE primitives, but does not touch arbitrary real-analytic primitives.  Is there a real-analytic counterexample outside the LE world, or can one prove a broader real-analytic impossibility theorem?
\end{itemize}
A third direction, not pursued here, is complexity.  Odrzywo\l{}ek tabulates EML tree sizes empirically but gives no lower bounds \cite{Odrzywolek2026}.  The Hardy-field and exp-height invariants suggested by the present work may provide a route to the first nontrivial lower bounds for EML depth.

\subsection*{Scope relative to Liouville--Ritt and Schanuel}
Our results concern single binary generators and Odrzywo\l{}ek's concrete calculator basis, not the full differential-algebraic notion of elementary extension.  In particular, Section~\ref{sec:eml} does not prove that every element of an arbitrary Liouville tower with algebraic adjunctions is representable by $\Edown$; the algebraic step beyond radicals remains outside the scope of this note.  Conversely, Theorem~\ref{thm:diagonal} is broader than any exp-log normal-form statement: it applies to every holomorphic or real-analytic binary germ with constant diagonal, whether elementary or not, and Theorem~\ref{thm:hardy} applies to the entire real logarithmico-exponential class through Hardy-field methods.

It is also worth separating the present arguments from Schanuel-type issues.  Odrzywo\l{}ek invokes Schanuel only to motivate a numerical sieve for the discovery phase \cite{Odrzywolek2026}.  None of the theorems proved here depends on Schanuel's conjecture.  The diagonal and Hardy-field obstructions are unconditional, and the EML section is a branch-correct symbolic core plus an honest accounting of the present computer-assisted evidence.  For related discussion of exp-log fields, closed forms, and the limited role of Schanuel in this neighborhood, see Chow \cite{Chow1999}.

\section*{Acknowledgment}
This manuscript is based on a close reading of Odrzywo\l{}ek's paper and public verification code, together with subsequent independent mathematical analysis.

\begin{thebibliography}{99}

\bibitem{AschenbrennerVanDenDriesVanDerHoeven2017}
Matthias Aschenbrenner, Lou van den Dries, and Joris van der Hoeven,
\emph{Asymptotic Differential Algebra and Model Theory of Transseries},
Annals of Mathematics Studies, vol.~195, Princeton University Press,
Princeton, NJ, 2017.

\bibitem{Chow1999}
Timothy Y. Chow,
\emph{What is a closed-form number?},
Amer. Math. Monthly \textbf{106} (1999), no.~5, 440--448.

\bibitem{Evans1989}
Trevor Evans,
\emph{Embedding and representation theorems for clones and varieties},
Bull. Austral. Math. Soc. \textbf{40} (1989), 199--205.

\bibitem{Goldstern2010}
Martin Goldstern,
\emph{A single binary function is enough},
unpublished note, 2010,
\url{https://dmg.tuwien.ac.at/goldstern/www/papers/notes/singlebinary.pdf}.

\bibitem{GoldsternPinsker2008}
Martin Goldstern and Michael Pinsker,
\emph{A survey of clones on infinite sets},
Algebra Universalis \textbf{59} (2008), 365--403.

\bibitem{Hardy1910}
G. H. Hardy,
\emph{Orders of Infinity: The `Infinit\"arcalcul' of Paul du Bois-Reymond},
Cambridge University Press, Cambridge, 1910.

\bibitem{Odrzywolek2026}
Andrzej Odrzywo\l{}ek,
\emph{All elementary functions from a single operator},
arXiv:2603.21852v2, 2026.

\bibitem{OdrzywolekRepo2026}
Andrzej Odrzywo\l{}ek,
\emph{SymbolicRegressionPackage},
GitHub repository, 2026,
\url{https://github.com/VA00/SymbolicRegressionPackage}.

\bibitem{OdrzywolekVerifier2026}
Andrzej Odrzywo\l{}ek,
\emph{verify\_eml\_symbolic\_chain.wl},
file in the repository \emph{SymbolicRegressionPackage}, 2026,
\url{https://github.com/VA00/SymbolicRegressionPackage/blob/master/EML_toolkit/EmL_verification/verify_eml_symbolic_chain.wl}.

\bibitem{Risch1969}
Robert H. Risch,
\emph{The problem of integration in finite terms},
Trans. Amer. Math. Soc. \textbf{139} (1969), 167--189.

\bibitem{Ritt1948}
J. F. Ritt,
\emph{Integration in Finite Terms: Liouville's Theory of Elementary Methods},
Columbia University Press, New York, 1948.

\bibitem{Rosenlicht1983}
Maxwell Rosenlicht,
\emph{Hardy fields},
J. Math. Anal. Appl. \textbf{93} (1983), no.~2, 297--311.

\bibitem{Sheffer1913}
Henry M. Sheffer,
\emph{A set of five independent postulates for Boolean algebras, with application to logical constants},
Trans. Amer. Math. Soc. \textbf{14} (1913), no.~4, 481--488.

\bibitem{Strzebonski2019}
Adam Strzebo\'nski,
\emph{Asymptotic solutions of polynomial equations with exp-log coefficients},
arXiv:1904.06796, 2019.

\end{thebibliography}

\end{document}
